\documentclass[a4paper,12pt,twoside,openright,titlepage]{book}

%Additional packages
\input{styles/packages}

% Listing style
\input{styles/lstset}

% Uncomment for production
% \syntaxonly

% Style
\pagestyle{headings}

% Turn on indexing
\makeindex[intoc]

% Define document
\author{D. Leeuw}
\title{Linux: Systeem Inventarisatie}
\date{\today\\v.0.9.5}

\begin{document}
\selectlanguage{dutch}

\maketitle

\copyright\ 2020-2025 Dennis Leeuw\\

\input{styles/licentie}

%%%%%%%%%%%%%%%%%%%
%%% Introductie %%%
%%%%%%%%%%%%%%%%%%%

\frontmatter
\chapter{Over dit Document}
\section{Leerdoelen}
\input{src/leerdoelen_cli}
\section{Voorkennis}
\input{src/voorkennis_cli}

%%%%%%%%%%%%%%%%%
%%% De inhoud %%%
%%%%%%%%%%%%%%%%%
\tableofcontents

\mainmatter

% Requires: sudo, cat
% Provides: lscpu, free, dmidecode, uname
% lspci, lsusb, lsscsi, lsblk, du, df, fdisk

\chapter{Systeem inventarisatie}
\input{src/inventarisatie}

\section{Systeem versienummers}
\input{src/versienummers}
\subsection{Distributie versie}
\input{src/distro_version}
\subsection{Kernel versie}
\input{src/kernel_version}

\section{Het moederbord}
\subsection{CPU}\index{CPU}\index{Hardware informatie!CPU}
\input{src/cpu_info}
\subsection{RAM}\index{RAM}\index{Hardware informatie!RAM}
Het geheugengebruik op een machine is belangrijk om te weten. Er zijn dan ook verschillende commando's die hier informatie over geven. In deze paragraaf zullen we het hebben over de commando's die je vertellen hoeveel geheugen er in de computer zit. Het belangrijkste is doe de Linux-kernel tegen het geheugen aankijkt en dat kunnen we zie door gebruikt te maken van \texttt{/proc/meminfo}\index{/proc/meminfo} bestand. Met \texttt{cat} of \texttt{less} kunnen we dat bestand lezen en dan zien we heel veel informatie over het geheugen. De belangrijkst informatie voor ons op dit moment is de totale hoeveelheid geheugen (RAM) die er in de machine zit en hoeveel er in gebruik is. Een compactere manier om dit weer te geven is door gebruik te maken van het \texttt{free}\index{free}\index{commando!free} commando:
\begin{lstlisting}[language=bash]
$ free 
         total      used      free   shared buff/cache   available
Mem:  32538364  18631520    591412  2071772   13315432    11378676
Swap: 15236492     63816  15172676
\end{lstlisting}
\texttt{free} heeft nog een andere handige functie en dat is dat het deze informatie met een bepaalde interval, bijvoorbeeld elke 10 seconden, kan weergeven. Zo kan je een drukke server in de gaten houden. Om dit te doen gebruik je \texttt{free -s 10}.

\subsection{BIOS}\index{BIOS}\index{Hardware informatie!BIOS}
De BIOS of EUFI is verantwoordelijk voor het opstarten van de computer, het moet dus heel veel weten van de hardware in een systeem. Als we de BIOS zouden kunnen uitvragen over ons systeem dan zou dat heel veel informatie kunnen opleveren, potentieel zelfs heel gevoelige informatie. Je kan dit dan ook niet als gewone gebruiker. Je moet root zijn om \texttt{dmidecode}\index{demidecode}\index{commando!demidecode} te kunnen gebruiken. In de gegeven voorbeelden zullen we dna ook \texttt{sudo} gebruiken.

De totale informatie die aanwezig is in de DMI-tabel is heel veel, dus het is handiger om er alleen de stukjes uit te halen die we nodig hebben. Dit doen we door aan \texttt{demidecode} de optie \texttt{-s} mee te geven met het onderdeel dat we zouden willen hebben. Als we bijvoorbeeld alleen de informatie over de BIOS versie willen hebben gebruiken we:
\begin{lstlisting}[language=bash]
$ sudo dmidecode -s bios-version
GNET88WW (2.36 )
\end{lstlisting}
Als je alleen de \texttt{-s} optie meegeeft en verder niets dan zie je welke keywords er bestaan voor \texttt{-s}.

Als we iets meer willen weten van het BIOS kunnen we ook meer informatie op vragen. Alle informatie over het systeem is opgedeeld in types. Met \texttt{-t} kan je alle informatie van een type opvragen, zonder een keyword krijg je een lijst te zien van beschikbare types. Alle BIOS informatie valt onder het keyword bios.
\begin{lstlisting}[language=bash]
$ sudo dmidecode -t bios
# dmidecode 3.2
Getting SMBIOS data from sysfs.
SMBIOS 2.7 present.

Handle 0x002A, DMI type 13, 22 bytes
BIOS Language Information
	Language Description Format: Abbreviated
	Installable Languages: 1
		en-US
	Currently Installed Language: en-US

Handle 0x0040, DMI type 0, 24 bytes
BIOS Information
	Vendor: LENOVO
	Version: GNET88WW (2.36 )
	Release Date: 05/30/2018
	Address: 0xE0000
	Runtime Size: 128 kB
	ROM Size: 12288 kB
	Characteristics:
		PCI is supported
		PNP is supported
		BIOS is upgradeable
		BIOS shadowing is allowed
		Boot from CD is supported
		Selectable boot is supported
		ACPI is supported
		USB legacy is supported
		BIOS boot specification is supported
		Targeted content distribution is supported
		UEFI is supported
	BIOS Revision: 2.36
	Firmware Revision: 1.14
\end{lstlisting}
Je ziet dat we zo al heel veel meer informatie naar boven halen.

\texttt{dmidecode} commando waarmee je veel informatie kan vinden. Lees de manual-pagina eens door en speel eens met de \texttt{-t} en \texttt{-s} opties om ermee vertrouwd te raken.




\section{Extensie bussen}
\subsection{PCI}\index{PCI}\index{Hardware informatie!PCI}
\input{src/pci_info}
\subsection{USB}\index{USB}\index{Hardware informatie!USB}
\input{src/usb_info}
\subsection{SCSI}\index{SCSI}\index{Hardware informatie!SCSI}
\input{src/scsi_info}

\section{Harddisks}
\input{src/disks}
\input{src/disks-fs}
Linux kent verscheidene tools om een harddisk te beheren, \'e\'en van de oudste en meest gebruikte is \texttt{fdisk}\index{fdisk}\index{commando!fdisk}. \texttt{fdisk} heeft geen grafische interface, maar wel een character gebaseerde interface. Naast \texttt{fdisk} zijn er ook \texttt{cfdisk}\index{cfdisk}\index{commando!cfdisk}, \texttt{sfdisk}\index{sfdisk}\index{commando!sfdisk} en \texttt{parted}\index{parted}\index{commando!parted}.


\input{src/disk-free}
\input{src/disk-usage}

\section{Netwerk}
\input{src/netwerk_hardware}
\subsection{Hostname}
De hostname van een GNU/Linux systeem is terug te vinden in het \texttt{/etc/hostname} bestand.



\section{De /etc directory}
\subsection{Datum en tijd}
De wereld is opgedeeld in tijdzones\index{tijdzone} en daarnaast heeft een land vaak nog regels met betrekking tot zomer- en wintertijd. De bestanden die de regels bevatten over wat de regels zijn per land of gebied zijn te vinden in \texttt{/usr/share/zoneinfo/}\index{zoneinfo}. Omdat de afspraak is dat configuratie bestanden in \texttt{/etc} staan is er een link vanuit \texttt{/etc/localtime}\index{localtime} naar de tijdzone die voor deze computer geldt in \texttt{/usr/share/zoneinfo/}.

Om de tijdzone aan te passen gebruiken we het commando:
\begin{lstlisting}
sudo dpkg-reconfigure tzdata
\end{lstlisting}\index{tzdata}


\subsection{Taal}
Het bestand \texttt{/etc/locale} bevat de gegevens over de taal die gebruikt wordt op het systeem. Een gebruiker kan in zijn \texttt{profile} een andere taal zetten als hij dat wil en deze door het systeem ondersteund wordt.

Het wijzigen van de locale, kan op een Debian systeem met:
\begin{lstlisting}[language=bash]
$ sudo dpkg-reconfigure locales
\end{lstlisting}
Eerst moet je een keuze maken welke talen er beschikbaar zijn op het systeem. Vink deze aan door gebruik te maken van de spatiebalk. Daarna selecteer je de default taal. Als je klaar bent reboot je het systeem om de wijzigingen te activeren.


\subsection{Keyboard}
Op een Debian systeem kan je de keyboard-layout wijzigen met de volgende commando's:
\begin{lstlisting}[language=bash]
$ sudo dpkg-reconfigure keyboard-configuration
$ sudo update-initramfs -u
\end{lstlisting}
Na de wijziging is er een reboot nodig om de wijzigingen te activeren.



\subsection{Login berichten}
\subsubsection{Issue}
In de \texttt{/etc/} directory vind je een bestand met de naam \texttt{issue}. Dit bestand bevat tekst die op het scherm afgebeeld wordt voor de Login prompt op de console. Standaard bevat het de naam en het versienummer van het systeem. In \texttt{issue} kan je die tekst wijzigen.


\subsubsection{motd}
In de \texttt{/etc/} directory vind je een bestand met de naam \texttt{motd}. MOTD staat voor Message Of The Day, ofwel het bericht van de dag. Dit geeft de systeembeheerder de mogelijkheid om belangrijke zaken met gebruikers te delen die op de commandline. Standaard bevat dit bestand een tekst over copyrights en de warranty.



\section{Inventarisatie opdracht}
\input{src/inventarisatie_opdracht}

%%%%%%%%%%%%%%%%%%%%%
%%% Index and End %%%
%%%%%%%%%%%%%%%%%%%%%
\backmatter
\printindex
\end{document}

%%% Last line %%%
